\documentclass[11pt,a4paper]{article}

% Packages
\usepackage[utf8]{inputenc}
\usepackage[T1]{fontenc}
\usepackage[italian]{babel}
\usepackage{geometry}
\usepackage{setspace}
\usepackage{enumitem}

\geometry{margin=2.5cm}
\setstretch{1.15}

\begin{document}

    \begin{center}
        {\Large \textbf{Documento dei Requisiti}} \\
        \vspace{0.5cm}
        \textbf{Nome del Progetto:} NutriApp \\
        \textbf{Data:} 28/10/2025
    \end{center}

    \section*{Informazioni Generali}

    \begin{itemize}
        \item \textbf{Committente:} Mileo Manuela
        \item \textbf{Contatti:} manuela.mileo@galileo.galileicrema.it
        \item \textbf{Autori:}

        \begin{itemize}
            \item Barakat Ismail
            \item Donzelli Christian
            \item Maltese Emanuel
            \item Yu Serena
        \end{itemize}

    \end{itemize}

% ======================================================

    \section*{Requisiti di Contesto}

    Il progetto NutriApp nasce per offrire uno strumento digitale avanzato che consenta agli utenti di monitorare e migliorare la propria alimentazione.

    È rivolto a chi desidera controllare i nutrienti assunti durante i pasti o necessita di un monitoraggio specifico per esigenze di salute (es. sportivi, diabetici, persone con intolleranze alimentari).

% ======================================================

    \section*{Requisiti degli Stakeholders}

    \begin{itemize}
        \item Monitoraggio completo dei valori nutrizionali (macro e micro nutrienti e acqua)

        \item Allarmi intelligenti personalizzabili in caso di superamento o avvicinamento alle soglie

        \item Tracciamento opzionale dei parametri fisici (peso, pressione, glicemia, ecc.) con analisi nel tempo

        \item Visualizzazione grafica dei dati per una comprensione immediata

        \item Personalizzazione degli obiettivi e dei parametri nutrizionali in base alle esigenze individuali

    \end{itemize}

% ======================================================

    \section*{Requisiti Funzionali}

    \subsection*{MUST}

    \begin{itemize}
        \item Monitoraggio completo dei valori nutrizionali

        \item Allarmi intelligenti (nel caso di superamento o avvicinamento alla soglia)

    \end{itemize}

    \subsection*{SHOULD}

    \begin{itemize}
        \item Grafici
    \end{itemize}

    \subsection*{MAY}

    \begin{itemize}
        \item Tracciamento opzionale
    \end{itemize}

% ======================================================

    \section*{Requisiti Non Funzionali}

    Il sistema dovrà garantire:

    \begin{itemize}

        \item \textbf{Usabilità:} interfaccia intuitiva, accessibile e user-friendly

        \item \textbf{Sicurezza:} crittografia dei dati sensibili e protezione della privacy degli utenti

        \item \textbf{Affidabilità:} riduzione dei rischi di crash o perdita dei dati

        \item \textbf{Scalabilità:} capacità di supportare un numero crescente di utenti senza perdita di prestazioni

        \item \textbf{Portabilità:} compatibilità con più piattaforme (Android, iOS, Web)

    \end{itemize}

\end{document}
